Para uma matriz real \(n \times n\) \(A \in M_n(\mathbb{R})\), denotamos por \(A^{\mathsf{R}}\) sua rotação anti-horária de \(90°\). Por exemplo,
\[
\begin{bmatrix} 1 & 2 & 3 \\ 4 & 5 & 6 \\ 7 & 8 & 9 \end{bmatrix}^{\mathsf{R}} = \begin{bmatrix} 3 & 6 & 9 \\ 2 & 5 & 8 \\ 1 & 4 & 7 \end{bmatrix}.
\]
Prove que se \(A = A^{\mathsf{R}}\), então para qualquer autovalor \(\lambda\) de \(A\), temos \(\operatorname{Re} \lambda = 0\) ou \(\operatorname{Im} \lambda = 0\).
